% ============================================================
% CHAPTER 5 — Experimental Setup (REVISED)
% Único cambio: la sección "Baselines" estaba declarada dos veces
% consecutivas (mismo \section, mismo \label{sec:baselines} repetido --
% LaTeX habría dado un warning de label duplicado). Se ha eliminado la
% primera versión (más genérica, describía 4 grupos a alto nivel) y se
% ha conservado la segunda (más precisa, describe solo los métodos
% realmente ejecutados) junto con el párrafo "Prior approaches" que ya
% seguía después. Nada más se ha tocado.
% ============================================================

\chapter{Experimental Setup}
\label{chap:experimental_setup}

This chapter defines how the three contributions of
Chapter~\ref{chap:method} are evaluated. The evaluation is organised
around five evaluation questions (Section~\ref{sec:experimental_objectives}).
Section~\ref{sec:benchmarks} describes the two benchmark processes,
Section~\ref{sec:training_protocol} the training protocol,
Section~\ref{sec:baselines} the comparison methods, and
Section~\ref{sec:metrics} the metrics. Section~\ref{sec:reproducibility}
gives the information needed to replicate the experiments.

% ------------------------------------------------------------
\section{Evaluation Questions}
\label{sec:experimental_objectives}

The five evaluation questions (EQs) below turn the sub-research
questions of Section~\ref{sec:research_questions} into concrete,
measurable experiments. Table~\ref{tab:rq_traceability} maps each EQ to
the research question and contribution it serves, and to the chapter
where its evidence appears.

\begin{description}
    \item[EQ1 --- Correctness.]
    Can the agent recover the exact Pareto front on a process whose
    optimal solution is known from an exact solver? (Contribution~C2.)

    \item[EQ2 --- Compliance.]
    Does the shield prevent invalid process states throughout
    training, and how fast does the agent learn to avoid illegal
    actions? (Contribution~C1.)

    \item[EQ3 --- Scalability.]
    How does the framework behave on a larger graph with a more complex
    trace space? (Contributions~C2 and~C3.)

    \item[EQ4 --- Comparative positioning and added value.]
    Does the framework reproduce the cost--duration trade-offs reported
    in prior work, and what does role-conditioned, capacity-constrained
    allocation add beyond independent per-event role selection?
    (All three contributions, qualitatively.)

    \item[EQ5 --- Component contribution.]
    What does each component add on its own, and does combining \ac{PPO}
    with shielding beat either alone? (Contributions~C1 and~C2, via the
    ablation of Chapter~\ref{chap:internal_validation}.)
\end{description}

EQ1 to EQ4 are addressed in Chapters~\ref{chap:results}
and~\ref{chap:discussion}; EQ5 in
Chapter~\ref{chap:internal_validation}. All five are revisited in
Section~\ref{sec:eq_revisited}.

\begin{table}[H]
\centering
\renewcommand{\arraystretch}{1.6}
\setlength{\tabcolsep}{12pt}
\caption{Traceability from sub-research questions to evaluation
         questions, contributions, and evidence chapters.}
\label{tab:rq_traceability}
\begin{tabular}{llll}
\toprule
\textbf{Research question} & \textbf{EQs} &
\textbf{Contrib.} & \textbf{Evidence} \\
\midrule
SRQ1 --- \ac{DCR} representation & EQ2 & C1 &
Ch.~\ref{chap:method}, \ref{chap:results} \\
SRQ2 --- Pareto discovery & EQ1, EQ3, EQ5 & C2 &
Ch.~\ref{chap:internal_validation}, \ref{chap:results} \\
SRQ3 --- Reward design & EQ1, EQ2, EQ5 & C2 &
Ch.~\ref{chap:internal_validation}, \ref{chap:results} \\
SRQ4 --- Role-conditioned allocation & EQ3, EQ4 & C3 &
Ch.~\ref{chap:results}, \ref{chap:discussion} \\
MRQ --- Compliant optimisation & EQ1--EQ5 & C1--C3 &
\S\ref{sec:eq_revisited} \\
\bottomrule
\end{tabular}
\end{table}

% TODO (Ch4): alinear la definición de C3 con esta versión honesta.
% C3 ya no es "resource allocation" sin más, sino "optimización
% role-conditioned del coste/duración, extendida a allocation bajo
% capacidad limitada del rol Expert". Revisar el enunciado de C3 en
% Methodology para que coincida con SRQ4 y con la sección de resultados.

% ------------------------------------------------------------
\section{Benchmarks}
\label{sec:benchmarks}
The framework is evaluated on two \ac{DCR} processes chosen for
complementary reasons. The Expense Report is small and has a known
exact solution, which makes it suitable for verifying correctness: the
agent's output can be checked against a provably optimal reference. The
Loan Application is larger, role-aware, and drawn from the same domain
as prior work, which makes it suitable for testing scalability and for
positioning the framework against existing approaches. Together they
cover both ends of the evaluation: a controlled correctness check on a
simple process, and a realistic stress test on a complex one. The two
benchmarks are described in turn below.

\subsection{Expense Report}
\label{sec:expense_dataset}

The first benchmark, used for EQ1 and EQ2, is the Expense Report
\ac{DCR} graph of Diaz et
al.~\cite{Diaz2024Pareto-OptimalModels}. It was chosen because an exact
Pareto front is available from a CSP+COP solver, allowing a direct
correctness check against the \ac{RL} agent. The process has five
activities and two objectives, cost and duration
(Table~\ref{tab:expense_costs}, Figure~\ref{fig:expense_graph}). With
five events and no roles, the observation space has $5 \times 3 = 15$
dimensions and the action space has 5 actions.

\begin{table}[H]
\centering
\renewcommand{\arraystretch}{1.6}
\setlength{\tabcolsep}{16pt}
\caption{Expense Report event costs and durations.}
\label{tab:expense_costs}
\begin{tabular}{lrr}
\toprule
\textbf{Event} & \textbf{Cost} & \textbf{Duration (min)} \\
\midrule
FillOutExpenseReport  & 10  & 1{,}500 \\
Approve               & 25  & 7{,}200 \\
Reject                & 30  & 7{,}200 \\
PayOut                & 15  & 1{,}800 \\
WithdrawExpenseReport & 700 & 1{,}200 \\
\bottomrule
\end{tabular}
\end{table}

\begin{figure}[H]
    \centering
    \includegraphics[width=0.7\linewidth]{Pictures/ExpenseGraph_DCR.png}
    \caption{Expense Report \ac{DCR} graph
    \cite{Diaz2024Pareto-OptimalModels}.}
    \label{fig:expense_graph}
\end{figure}

\subsection{Loan Application}
\label{sec:loan_dataset}

The second benchmark, used for EQ3 and EQ4, is the Loan Application
process, the LoanJS variant studied by Kirchdorfer et
al.~\cite{KirchdorferLukasFromPolicies}. Its event log is generated by
simulation with AgentSimulator and contains 1{,}000 cases over 12
activities, obtained from the authors' public repository. This
benchmark was chosen because Kirchdorfer et al.\ optimise the same
process with an NSGA-II approach, enabling the comparison of EQ4.

\paragraph{Graph construction.}
The graph structure was mined from the log: the log was imported into Disco~\cite{Fluxicon2012Disco:Software}, whose Fuzzy Miner produced the control-flow model shown in Appendix~\ref{app:loanapp_disco}. This model was then imported into DCR-js, where conditions, responses, and exclusions were added to reflect the domain logic. The result is a \ac{DCR} graph with 12 events and 61 relations. Because the full graph is too dense to read at page width, it is reproduced, together with a complete relation table, in Figure \ref{fig:loanapp_dcrjs_full}.


\begin{figure}[H]
    \centering
    \includegraphics[width=0.6\textheight]{Pictures/LoanAppDCR.png}
    \caption{Loan Application \ac{DCR} graph, modelled in DCR-js from
             the Disco-discovered structure of
             Figure~\ref{fig:loanapp_disco}. Arcs encode conditions,
             responses, and exclusions; each event additionally carries
             the global Junior/Expert role multipliers of
             Section~\ref{sec:roles}.}
    \label{fig:loanapp_dcrjs_full}
\end{figure}


\paragraph{Costs and durations.}
The base durations come from the log: each value is the mean duration
of that activity across the 1{,}000 cases, rounded to the nearest
minute. The base costs were assigned by hand, because the log records
no cost data. Kirchdorfer et al.\ faced the same gap and likewise added
a synthetic cost model~\cite{KirchdorferLukasFromPolicies}. Costs were
set to reflect the relative effort of each activity, so the
expert-review steps (\emph{Approve loan offer}, \emph{AML check}) carry
the highest cost. Table~\ref{tab:loan_costs} lists the values.


\begin{table}[H]
\centering
\renewcommand{\arraystretch}{1.5}
\setlength{\tabcolsep}{16pt}
\caption{Loan Application base costs and durations. Durations are the
         rounded per-activity means of the event log; costs are
         hand-assigned. Per-role values follow the multipliers of
         Section~\ref{sec:roles}.}
\label{tab:loan_costs}
\begin{tabular}{lrr}
\toprule
\textbf{Event} & \textbf{Cost} & \textbf{Duration (min)} \\
\midrule
Check application form completeness &  76  &  55  \\
Appraise property                   & 509  & 244  \\
Check credit history                & 118  &  84  \\
AML check                           & 539  & 258  \\
Assess loan risk                    & 318  & 175  \\
Design loan offer                   &  35  &  19  \\
Approve loan offer                  & 767  & 184  \\
Approve application                 &  20  &  15  \\
Reject application                  &  40  &  29  \\
Cancel application                  &  20  &  14  \\
Return application to applicant     &  21  &  15  \\
Applicant completes form            &   1  & 478  \\
\bottomrule
\end{tabular}
\end{table}

\paragraph{Roles and action space.}
Each event can be executed by a \emph{Junior} or an \emph{Expert}
resource, using the multipliers of Section~\ref{sec:roles}. This
follows the Junior/Senior structure of Kirchdorfer et
al.~\cite{KirchdorferLukasFromPolicies}, supporting the comparison of
EQ4. The action space therefore has $12 \times 2 = 24$ event--role
pairs, and the observation space has $12 \times 3 = 36$ dimensions. No
exact Pareto front exists for this graph, so it is evaluated through
convergence behaviour, Pareto coverage, and role-selection patterns.

\paragraph{Capacity-constrained variant.}
Under the role model above, both roles are always available, so the
weight-optimal role of each event can be chosen independently of the
others: role selection is separable across events and does not yet
constitute resource allocation. To introduce the scarcity that makes
allocation a non-trivial decision, a capacity-constrained variant
limits the \emph{Expert} role to a fixed budget of $k$ assignments per
case. The remaining Expert budget is exposed to the agent as an
additional observation dimension, and the limit is enforced by the
\ac{DCR} shield: once the budget is exhausted, Expert actions are
treated as not enabled and are intercepted exactly as structurally
illegal events are (Chapter~\ref{chap:method}). A finite budget couples
the per-event role decisions, since spending the scarce Expert
allocation on one activity makes it unavailable for another, so the
agent must learn \emph{where} to spend it rather than deciding each
event in isolation. The unconstrained case $k=\infty$ recovers the
role-conditioned experiment described above and serves as its baseline.
The budget is swept over $k \in \{1, 2, 3, \infty\}$
% TODO: confirmar el conjunto de valores de k tras calibrar; elegir k
% que sean vinculantes, es decir, menores que el número de Expert que
% el agente usaría sin restricción bajo el peso duración-dominante.
to characterise how the recovered front and the Expert-allocation
pattern respond to scarcity. This variant captures opportunity-cost
allocation within a single case; it does not model inter-case
contention, queues, waiting time, or utilisation, and the comparison
with simulation-based approaches therefore remains structural rather
than numerical (Section~\ref{sec:loanapp_interpretation}).

\begin{remark}
A preliminary run without roles showed that 50{,}000 steps are not
enough to recover the full Pareto front on this graph: the agent found
only one of three accepting trace variants. This motivates the larger
budget used with roles and is discussed in
Section~\ref{sec:discussion_scalability}.
\end{remark}

% ------------------------------------------------------------
\section{Training Protocol}
\label{sec:training_protocol}

All experiments use the \ac{PPO} implementation of Stable
Baselines~3~\cite{Raffin2021Stable-Baselines3:Implementations}, with the hyperparameters of
Table~\ref{tab:hyperparams}. The Expense Report and the role-conditioned
Loan Application experiments train for 100{,}000 steps; the preliminary
no-roles run used 50{,}000. The capacity-constrained variant uses the
same 100{,}000-step protocol per value of $k$. An episode is cut off
after 300 steps if it has not reached an accepting state, and is then
recorded as non-accepting.

Six weight configurations are used to explore objective preferences:
\[
(\alpha,\beta) \in
\{(0,0),\,(1,0),\,(0,1),\,(0.5,0.5),\,(2,0.5),\,(0.5,2)\}.
\]
They span the trade-off from pure cost to pure duration, with balanced
and asymmetric points in between (see
Section~\ref{sec:scalarisation}). Each configuration is trained
separately. Following Section~\ref{sec:training_pipeline}, all reported
statistics come from the accepting traces of the final 20\,\% of each
run; there is no separate evaluation phase.

% ------------------------------------------------------------
%% CAMBIO: la sección "Baselines" estaba duplicada (dos \section, dos
%% \label{sec:baselines} idénticos). Se conserva una sola vez, con el
%% párrafo introductorio de la versión más precisa (la que describe la
%% tabla de métodos realmente ejecutados) seguido del párrafo de
%% "Prior approaches" que ya estaba a continuación.
\section{Baselines}
\label{sec:baselines}

Table~\ref{tab:implemented_baselines} lists the methods implemented and
executed in this thesis, with the evaluation question each one serves.

\begin{table}[H]
\centering
\renewcommand{\arraystretch}{1.35}
\caption{Baselines implemented and executed in this thesis, with the
         evaluation question each one serves.}
\label{tab:implemented_baselines}
\begin{tabular}{@{}p{3.2cm} p{1.0cm} p{7.2cm}@{}}
\toprule
\textbf{Method} & \textbf{EQ} & \textbf{Purpose} \\
\midrule
CSP+COP (MiniZinc) & EQ1 &
Exact Pareto front used as ground truth on the Expense Report
benchmark~\cite{Diaz2024Pareto-OptimalModels}. \\
Unshielded \ac{PPO} & EQ5 &
Isolates the effect of removing the \ac{DCR} shield during learning. \\
Shielded random agent & EQ5 &
Isolates the effect of structural compliance without learning. \\
Always-Junior & EQ4 &
Cheapest, slowest valid role policy; bounds the cost axis of the
role-conditioned front. \\
Always-Expert & EQ4 &
Fastest, most expensive role policy; bounds the duration axis in the
unconstrained case, where the Expert budget is unlimited. \\
Greedy per-event & EQ4 &
Assigns each event its weight-optimal role independently; the separable
optimum under unlimited capacity, and the naive heuristic the learned
policy must beat under a finite budget. \\
\bottomrule
\end{tabular}
\end{table}

\paragraph{Prior approaches (qualitative positioning).}
The framework is additionally positioned against three prior approaches:
the \ac{RL}+\ac{BPM} method of Huang et
al.~\cite{Huang2011ReinforcementManagement}, the NSGA-II
handover-policy optimisation of Kirchdorfer et
al.~\cite{KirchdorferLukasFromPolicies}, and the simulation-based trace
generation of Doumeni~\cite{ArtemisDoumeniParetoAgentSimulator}. These
are not baselines in the experimental sense: they are not re-executed on
the present implementation, and their reported figures are not directly
comparable, since they use different cost models, metrics (labour cost,
waiting time) and granularity (pool- or policy-level rather than
event-level). They are therefore used for qualitative positioning only.
The comparison is developed in Chapter~\ref{chap:discussion} under EQ4,
and each approach is introduced in full in
Chapter~\ref{chap:background}.
% TODO: ajustar \ref{chap:related_work} a tu label real.
% ------------------------------------------------------------
\section{Evaluation Metrics}
\label{sec:metrics}

Five metrics are used across the experiments:

\begin{description}
    \item[Acceptance rate.] Share of episodes ending in an accepting
    marking (no event both included and pending). EQ1, EQ2.

    \item[Illegal-action ratio.] Share of action attempts that violate
    the enablement conditions of the current marking; a falling ratio
    shows the agent learning legal behaviour. EQ2.

    \item[Pareto coverage.] Whether the accepting episodes contain all
    known Pareto-optimal traces (against the CSP+COP ground truth).
    EQ1, EQ3.

    \item[Convergence behaviour.] Training stability, read from the
    reward and illegal-action curves. EQ2, EQ3.

    \item[Role-selection distribution.] Share of Junior and Expert
    choices per event across weight configurations and, in the
    capacity-constrained variant, the distribution of the scarce Expert
    budget across events. Used to assess whether the agent concentrates
    Expert effort on the events where it most affects the trade-off.
    EQ4.
\end{description}

% ------------------------------------------------------------
\section{Reproducibility}
\label{sec:reproducibility}

All code, process models, and notebooks are available at
\url{https://github.com/Sofiaortizarce/dcr-js} (branch
\texttt{feature/cost-roles}).

The \ac{DCR} graphs are under
\texttt{app/public/examples/diagrams/}: \texttt{Expense\_Report\_Diaz.xml},
\texttt{Loan\_Application\_Diaz.xml} (no roles), and
\texttt{LoanApp\_junior\_senior.xml} (with roles). Experiments are
launched via \texttt{scripts/run\_experiments.py}, which connects the
Node.js adapter, the Gymnasium environment, and the \ac{PPO} agent.
% TODO: fijar la versión de stable-baselines3 en requirements.txt.

All runs use a fixed seed $s = 1$. Results may vary marginally across
hardware due to GPU floating-point non-determinism; the sensitivity to
seed variation is evaluated in
Section~\ref{sec:policy_generalisation}, and the qualitative findings
(Pareto front shape, role specialisation) are stable across seeds. Raw
logs (about 500\,MB) are not stored in the repository but can be
regenerated from the provided configurations; the analysis notebooks
reproduce every figure and table of Chapter~\ref{chap:results}.
